% Generated draft: requires mathematical and source review.
\documentclass[11pt]{article}
\usepackage[margin=25mm]{geometry}
\usepackage{fontspec}
\setmainfont{Latin Modern Roman}
\usepackage{amsmath,amssymb,mathtools,unicode-math}
\setmathfont{Latin Modern Math}
\usepackage{enumitem,needspace,xurl,hyperref}
\setlength{\parindent}{0pt}
\setlength{\parskip}{5pt}
\setlength{\emergencystretch}{3em}
\newcommand{\N}{\mathbb{N}}
\newcommand{\Z}{\mathbb{Z}}
\newcommand{\Q}{\mathbb{Q}}
\newcommand{\R}{\mathbb{R}}
\newcommand{\C}{\mathbb{C}}
\newcommand{\F}{\mathbb{F}}
\newcommand{\A}{\mathbb{A}}
\newcommand{\PP}{\mathbb{P}}
\newcommand{\E}{\mathbb{E}}
\newcommand{\eps}{\varepsilon}
\newcommand{\dd}{\,\mathrm{d}}
\begin{document}
\section*{UnsolvedMath ID: 20002491}
\textbf{Source problem code:} catalogue entry
\subsection{Definitions and mathematical statement}
Diogo Gomes's Problem 35 concerns small perturbations of a completely integrable Hamiltonian and a nonresonant unperturbed torus. It reports uniform formal approximations of the corresponding viscosity solution and formal density expansions, and asks how well the Mather measures themselves are approximated. In a usual action-angle model, \(H_0=H_0(p)\) and nonresonance of a torus means \(k\cdot\nabla H_0(p_0)\ne0\) for all \(0\ne k\in\Z^n\). The source does not specify a perturbation family, regularity, Diophantine constants, norms, uniqueness of the measure, or a quantitative order to prove.

A Mather measure is an invariant probability minimizing the cohomology-adjusted Lagrangian action; its cotangent version is obtained by the Legendre map. The aid's transfer lemma is purely metric and applies whenever the relevant phase-space probabilities have graph representations. On a compact base \(X\) with reference measure \(\dd x\), let
\[
 \mu=(x\mapsto(x,G(x)))_*(m(x)\,\dd x),\qquad
 \widetilde\mu=(x\mapsto(x,\widetilde G(x)))_*(\widetilde m(x)\,\dd x),
\]
where \(m,\widetilde m\ge0\), \(\int m=\int\widetilde m=1\), and both graphs lie in a common phase-space metric set of finite diameter \(D\). Assume that the distance between two points over the same base point is bounded by the indicated fiber distance \(|G(x)-\widetilde G(x)|\), as in the product metric. Define \(W_1\) as the infimum, over probability couplings, of the mean phase-space distance. The aid gives
\[
 W_1(\mu,\widetilde\mu)
 \le\int_X\min\{m,\widetilde m\}|G-\widetilde G|\,\dd x
 +\frac D2\|m-\widetilde m\|_{L^1}.
\]
The common-density mass is matched over each base point and the remaining mass is half the \(L^1\) density difference.

Thus, for a perturbation parameter \(\varepsilon\) and order-\(N\) approximants defining genuine probabilities in one common bounded phase-space region, bounds
\[
 \|G_\varepsilon-G_{\varepsilon,N}\|_\infty
 \le C_G|\varepsilon|^{N+1},\qquad
 \|m_\varepsilon-m_{\varepsilon,N}\|_1
 \le C_m|\varepsilon|^{N+1}
\]
imply
\[
 W_1(\mu_\varepsilon,\mu_{\varepsilon,N})
 \le(C_G+DC_m/2)|\varepsilon|^{N+1}.
\]
A formal density truncation that is negative or unnormalized is not automatically a probability; any positivity or normalization correction needs its own error bound. For cotangent graphs \(G_\varepsilon=c+D u_\varepsilon\), the required estimate is on gradients, not just potentials. Conversion to velocity graphs requires corresponding control of the Legendre map.

The aid's smooth oscillatory obstruction can be made explicit on \(X=\R/\Z\):
\[
 u_j(x)=\frac{\sin(2\pi jx)}{2\pi j},\quad
 u_0=0,\quad m_j=m_0=1,\quad
 G_j=u_j'=\cos(2\pi jx),\quad G_0=0.
\]
Although \(\|u_j-u_0\|_\infty\to0\), in the product metric on \(X\times[-1,1]\) the graph probabilities satisfy
\[
 W_1(\mu_j,\mu_0)=\int_0^1|\cos(2\pi jx)|\,\dd x=2/\pi.
\]
This demonstrates the failure of a general graph-measure estimate based only on uniform potentials and exact density agreement. It is not asserted that these graphs are Mather measures for the specific near-integrable family in the original problem. The transfer lemma reduces that problem to invariant-graph and density remainder estimates; it does not itself prove those dynamical estimates or their small-divisor assumptions.
\subsection{Short English statement}
To approximate phase-space Mather measures, the aid bounds Wasserstein error by graph displacement and density error. Uniform approximation of the viscosity potential alone is insufficient: an oscillatory example has vanishing potential error but a fixed measure error. The needed invariant-graph and density remainder estimates remain separate dynamical work.
\subsection{Sources}
\begin{itemize}
\item Local registry, permanent UnsolvedMath ID \(20002491\), exact public code \texttt{catalogue entry}; canonical location \texttt{aim-pdes-notes.json}, \texttt{notes[102]}.
\item Diogo Gomes, Problem 35, AIM \emph{New connections between dynamical systems and PDE's}; supplied reference \url{https://aimath.org/WWN/dynpde/dynpde.pdf}. The verification field reports that the official workshop page corrects OCR number 5 to 35 and removes the isolated trailing page digit 3. No source was accessed here.
\item Supplied research attempt 1: \texttt{partial\_result}, run and registry status \texttt{partially\_solved}, rigor label \texttt{proved}, low novelty confidence; classification \texttt{PARTIAL-PROGRESS}, statement status \texttt{reconstructed\_unverified}, created and updated 14 August 2026. No proposer year, original remarks, literature entry, or primary citation for the coupling lemma is supplied.
\item The research aid explicitly says it is machine generated, not peer reviewed, and requires independent verification before citation.
\end{itemize}
\end{document}
